El cáncer colorrectal (CCR) es un problema de salud importante a nivel mundial, con una incidencia particularmente alta en Chile. Esta tesis se estructura en dos fases distintas: la primera tiene como objetivo entender y analizar el proceso de atención del CCR en el Servicio de Salud Metropolitano Sur (SSMS) del sistema de salud pública de Chile, mientras que la segunda se centra en ofrecer una estrategia de mejora factible para la atención del CCR en esta zona.

La fase de descubrimiento utilizó la Minería de Procesos, una disciplina innovadora que permite identificar rápidamente los casos más comunes y obtener indicadores de rendimiento para los pacientes del sistema. Este estudio utilizó datos reales del SSMS para analizar y evaluar la calidad de atención de CCR. El conjunto de datos, que abarca desde 2017 hasta 2023, incluye 2048 y 2524 casos de pacientes. A través de la colaboración con el Departamento de Gestión y Redes del SSMS, se revelaron el funcionamiento y los problemas que aquejan al servicio en el diagnóstico y tratamiento de CCR.

La fase de mejora presenta un modelo de microsimulación diseñado para simular la implementación de una estrategia de detección de CCR. Esta estrategia implica realizar pruebas inmunohistoquímicas fecales aleatorias a la población de entre 50 y 75 años. Quienes obtengan un resultado positivo en la prueba se deberán realizar una colonoscopia, con el objetivo de detectar la enfermedad en sus etapas iniciales. Todos los datos demográficos y médicos utilizados en el modelo se obtuvieron de fuentes respetadas, para garantizar la validez del modelo. Los resultados indican que a medida que aumentan las tasas de adherencia de la población al programa de detección, hay una reducción constante en los casos de cáncer en etapa avanzada; sin embargo, también aumentan los costos totales.

\vfill
\noindent {\bf Palabras Claves}: Salud, Minería de Procesos, Simulación, Optimización, Cáncer Colorrectal, Microsimulación, Tamizaje.
